\documentclass[10pt,a4paper]{amsart}
\usepackage[margin=19mm]{geometry}
\usepackage{amsmath,amssymb,mathtools}
\usepackage[scr=rsfs,cal=euler]{mathalfa}
\usepackage{xfrac}
\usepackage[unicode,hidelinks,bookmarks=false]{hyperref}
\newcommand{\ie}{i.e.\ }
\newcommand{\cf}{cf.\ }
\newcommand{\resp}{respectively\ }
\newcommand{\Subsection}{\subsection}
\providecommand{\nobreakdash}{\nobreak\hbox{-}}
\setlength{\emergencystretch}{2em}
\multlinegap0pt
\usepackage{tikz}
\usepackage{enumerate}
\usepackage{bm}
\usepackage{amssymb}
\usepackage{float}
\usepackage{xcolor}
\newtheorem{de}{Definition}
\newtheorem{lemme}{Lemma}
\newtheorem{prop}{Proposition}
\renewcommand{\geq}{\geqslant}
\renewcommand{\leq}{\leqslant}
\title{Homogeneity in Coxeter groups and split crystallographic groups}
\author{Simon Andr{\' e}, Gianluca Paolini}
\date{Source: arXiv:2504.18354v3 (CC BY 4.0)}
\begin{document}
\maketitle

\noindent\emph{Source and licence.} Homogeneity in Coxeter groups and split crystallographic groups. Version arXiv:2504.18354v3 (CC BY 4.0), distributed by arXiv under the Creative Commons Attribution 4.0 International licence, http://creativecommons.org/licenses/by/4.0/, as recorded in the arXiv licence metadata for this exact version and shown on the abstract page https://arxiv.org/abs/2504.18354v3. Full licence notice: \texttt{licenses/arxiv-2504.18354-CC-BY-4.0.txt}.\par\smallskip

\noindent\textit{Editorial scope.} Counted unit: the article's prop 15janvier2026 (source0.tex lines 421--423). The article's complete original proof of it is delivered with it (source0.tex lines 425--447, one \texttt{proof} environment of 23 lines). No mathematics is written, repaired or re-derived: the statement and the proof are the article's own lines, in the article's own notation, and no retained line is substituted or re-flowed.  Retained uncounted as the source-local context the proof needs: lines 303--303; lines 317--319; lines 374--374.  Nothing was omitted from any retained span, so the package carries no omission record.  No retained line contains a citation, so the wrapper carries no bibliography and no bibliography entry.  No retained line contains a figure environment, an image-inclusion command or drawing code, so no image is delivered and the package declares no asset dependency.  Editorial formatting only, in addition: an AMS \texttt{amsart} preamble that restates the article's own declarations and macros used by the retained lines, namely the environments \texttt{de}, \texttt{lemme}, \texttt{prop} on separate counters, together with the standard packages those lines need.  The retained environments print in this document as De 1, Lemme 1, Proposition 1 in order of appearance, whereas the article prints the same items with the numbering of its own full document.  The complete native source of the article as distributed in the arXiv e-print for this exact version is bundled unchanged as \texttt{sources/177089/source0.tex}.
\begin{de}\label{def_irreducible}We say that $G$ is \emph{irreducible} if $\rho$, viewed as a representation $G_0\rightarrow\mathrm{GL}_n(\mathbb{Q})$, is irreducible (meaning that the only linear subspaces of $\mathbb{Q}^n$ that are stable under the action of $\rho(G_0)$ are $\mathbb{Q}^n$ and $\lbrace 0\rbrace$).\end{de}

\begin{lemme}\label{finitely_many}
Let $G$ be a crystallographic group. Then $G$ has only finitely many conjugacy classes of finite subgroups.
\end{lemme}

\begin{de}\label{class-permuting}Let $G$ be a group and let $\varphi$ be an endomorphism of $G$. We say that $\varphi$ is \emph{class-permuting} if the following condition holds: for any two non-conjugate finite subgroups $H,H'$ of $G$, $\varphi(H)$ and $\varphi(H')$ are non-conjugate.\end{de}

\begin{prop}\label{15janvier2026}
Let $G=H\rtimes G_0$ be an irreducible split crystallographic group. Let $\ell\geq 1$ and $u=(u_1,\ldots,u_{\ell})$ be a tuple of elements of $G$. Suppose that the subgroup $U$ of $G$ generated by $\lbrace u_1,\ldots,u_{\ell}\rbrace$ is infinite and that there exists a class-permuting endomorphism $\theta$ of $G$ such that $\theta(u)=u$ and $\theta(H)\subset H$. Then $\theta$ is an automorphism of $G$.
\end{prop}

\begin{proof}
By Lemma \ref{finitely_many}, $G$ has only finitely many conjugacy classes of finite subgroups. Let $F_1,\ldots,F_m$ be a collection of representatives of these conjugacy classes. Let $\mathcal{C}=\lbrace [F_1],\ldots,[F_k]\rbrace$, where $[F_i]$ denotes the conjugacy class of $F_i$ in $G$. Since the set $\mathcal{C}$ is finite, and since $\theta$ permutes the conjugacy classes of finite subgroups of $G$ (see the remark following Definition \ref{class-permuting}), there is an integer $N\geq 1$ such that $\theta^N$ acts trivially on $\mathcal{C}$. Hence, $\theta^N([G_0])=[G_0]$ and thus there is an element $g\in G$ such that $\theta^N(G_0)=gG_0g^{-1}$. Write $g=hg_0$ for some $h\in H$ and $g_0\in G_0$. The endomorphism $\theta'=\mathrm{ad}(h^{-1})\circ \theta^N$ satisfies $\theta'(H)\subset H$ and $\theta'(G_0)=G_0$. As $G_0$ has finite automorphism group, there is an integer $M\geq 1$ such that $\theta'^M$ induces the identity of $G_0$. Define $f=\theta'^M$. As $\theta'=\mathrm{ad}(h^{-1})\circ \theta^N$, there is an element $h'\in H$ such that $\theta'^M=\mathrm{ad}(h')\circ \theta^{MN}$.

\smallskip \noindent
We will prove that $f$ is in fact the identity of $G$ (not only of $G_0$). It follows that $\theta'$ is an automorphism of $G$, and thus that $\theta$ is an automorphism of $G$.

%We will prove that $f$ is in fact the identity of $G$ (not only of $G_0$). But $f=\mathrm{ad}(h^{-1})\circ \theta^N$ with $\theta=\varphi'\circ \varphi$, so $\varphi'$ must be surjective, and thus $\varphi'$ must be an automorphism of $G$ (indeed $G$ is Hopfian, as it embeds in $\mathrm{GL}_{n+1}(\mathbb{Z})$). Moreover, $\varphi'(u')=u$, which shows that $u$ and $u'$ are automorphic in $G$.

\smallskip \noindent
It remains to prove that $f$ is the identity of $G$. Since $U$ is infinite (by assumption), it contains an element $x$ of infinite order (indeed, as well known, finitely generated linear periodic groups are finite), hence $y:=x^{\vert G_0\vert}$ is a non-trivial element of $H$. But $\theta(u)=u$, so $\theta_{\vert U}=\mathrm{id}_{ U}$ and $f_{\vert U\cap H}=\mathrm{id}_{U\cap H}$ (since $f=\mathrm{ad}(h')\circ \theta^{MN}$ with $h'\in H$), and therefore $f(y)=y$.

\smallskip \noindent
By assumption, $G$ is irreducible, which means that the linear representation $\rho : G_0\rightarrow \mathrm{GL}_n(\mathbb{Z})\subset \mathrm{GL}_n(\mathbb{Q})$ is irreducible (see Definition \ref{def_irreducible}), where $\rho$ denotes the action of $G_0$ on $H$ in the semidirect product $G=H\rtimes G_0$.  Let $V$ be the linear subspace of $\mathbb{Q}^n$ spanned by the finite set $G_0\cdot y$. By irreducibility and the fact that $y$ is non-trivial, $V$ must coincide with $\mathbb{Q}^n$. It follows that $G_0\cdot y$ contains $n$ vectors $h_1,\ldots, h_n\in H=\mathbb{Z}^n$ that are linearly independent over $\mathbb{Q}$. Let $H'$ be the subgroup of $H$ generated by $\lbrace h_1,\ldots, h_n\rbrace$ and let $A\in \mathrm{M}_n(\mathbb{Z})$ be the matrix whose columns are the vectors $h_1,\ldots,h_n$ written in the canonical basis $e_1=(1,0,\ldots,0),\ldots,e_n=(0,\ldots,0,1)$ of $H=\mathbb{Z}^n$. By the inverse of matrix formula, there is a matrix $B\in \mathrm{M}_n(\mathbb{Z})$ (namely $B=A^\intercal$) such that $AB=dI_n$, with $d=\mathrm{det}(A)\neq 0$ since $h_1,\ldots, h_n$ are linearly independent over $\mathbb{Q}$. It follows that $dH\subset H'$ (using additive notation).

%\emph{A fortiori}, these vectors are linearly independent over $\mathbb{Z}$, and thus $H'$ is a subgroup of $H=\mathbb{Z}^n$ isomorphic to $H$. 

%(in other words, the subgroup $H'=\langle h_1,\ldots, h_n\rangle$ is a lattice in $K^n$)

\smallskip \noindent
For each $1\leq i\leq n$, as $h_i$ belongs to $G_0\cdot y$, one can write $h_i=g_i(y)$ for some $g_i\in G_0$. Using the fact that $G_0$ acts on $H$ by conjugation in the semidirect product $H\rtimes G_0$, let us write $h_i=g_iyg_i^{-1}$. As we have proved in the previous paragraphs that $f(y)=y$ and that the restriction of $f$ to $G_0$ is the identity, we have $f(h_i)=f(g_iyg_i^{-1})=f(g_i)f(y)f(g_i)^{-1}=g_iyg_i^{-1}=h_i$ for each $1\leq i\leq n$. Hence $f$ coincides with the identity on $H'$.

\smallskip \noindent
Now, recall that $dH\subset H'=\langle h_1,\ldots,h_n\rangle$ with $d\in\mathbb{Z}^{\ast}$. Therefore, for every integer $1\leq i\leq n$, the element $de_i$ belongs to $H'$. But we have just proved that $f$ is the identity on $H'$, so we have $f(de_i)=de_i$, hence $df(e_i)=de_i$ and $f(e_i)=e_i$. Conclusion: $f_{\vert H}$ is the identity of $H$, and so $f$ is the identity of $G$.\end{proof}

\end{document}
